\documentclass[11pt]{article}
\usepackage{preamble}
\renewcommand{\answer}[1]{}

\begin{document}

\ladderheading{AP Statistics \textperiodcentered\ Topic 4.7 (CED 2.8) \textperiodcentered\ B: 2026-10-30 \textperiodcentered\ E: 2026-11-20}{Average Payoff, Typical Ticket}

\focusbanner{TODAY'S PROBLEM}

Suppose a club sells 200 raffle tickets for \$2 each. One ticket wins \$100, four tickets win \$25 each, 20 tickets win \$5 each, and the other tickets win nothing. Let $X$ be the buyer's net profit from one randomly chosen ticket: prize minus ticket cost.\par\smallskip \textbf{Priya:} \emph{The average gross payout is \$1.50 per ticket, so buying one is a good deal.}\par \textbf{Mateo:} \emph{Most tickets win nothing.}\par
\begin{center}
\textbf{Ticket net-profit distribution}\par\smallskip
\begin{tabular}{|l|c|c|}
\hline
\textbf{Prize} & \textbf{Net $x$} & \textbf{$P(X=x)$} \\ \hline
\textbf{Top prize} & $98$ & $1/200$ \\ \hline
\textbf{Second prize} & $23$ & $4/200$ \\ \hline
\textbf{Small prize} & $3$ & $20/200$ \\ \hline
\textbf{No prize} & $-2$ & $175/200$ \\ \hline
\end{tabular}
\end{center}
\begin{firsttakebox}{FIRST TAKE --- before the video (5 min)}
Would you buy one ticket if your goal were to make money? Use one number or outcome from the problem.
\workspace{0.9in}
\end{firsttakebox}

\begin{callout}{\IconBook\ DISTRIBUTION, EXPECTATION, AND ONE OUTCOME (keep on the page)}{calloutgreen}
A discrete random variable assigns a number to each random outcome. Its probability distribution lists every
possible value $x$ with $0\le P(X=x)\le1$, and the probabilities sum to 1. The expected value
$\mu_X=\sum xP(X=x)$ is the \textbf{long-run average} over many repetitions, not the result promised on one trial.
For a buyer's \textbf{net profit}, subtract the ticket cost from every gross prize before judging the deal.
\end{callout}

\newpage
\textbf{Finish after the video:}\par\smallskip

\dokbadge{1}~\textbf{(a)} State the four possible values of $X$ and identify the most likely value.\par
\workspace{0.7in}

\dokbadge{2}~\textbf{(b)} Verify that the probabilities form a distribution. Calculate the expected gross payout and the expected net profit $\mu_X$.\par
\workspace{1.4in}

\dokbadge{3}~\textbf{(c)} Decide whether Priya's numerical statement and Mateo's statement can both be true, and explain why there is no contradiction. Assess Priya's good-deal conclusion using the ticket cost and $\mu_X$. Explain what expected value does and does not predict for one buyer, then decide whether you would buy and defend the criterion behind your choice.\par
\workspace{2.0in}

\begin{sentenceframebox}\raggedright\textbf{Frame:} Both statements can be true because \blank[2.7in], while the large prizes \blank[1.7in].\end{sentenceframebox}

\begin{sentenceframebox}\raggedright\textbf{Frame:} I would \blank[0.7in] because $\mu_X=\blank[0.7in]$ means \blank[2.4in] in the long run, not \blank[1.8in] for one ticket.\end{sentenceframebox}

\begin{summaryexitbox}{EXIT --- turn this in}
One thing the video changed about my first take: \hrulefill\par
\workspace{0.4in}
\end{summaryexitbox}

\end{document}
